\documentclass[UTF8,zihao=-4]{ctexart}
\usepackage[a4paper,margin=2.2cm]{geometry}
\usepackage{xcolor}
\usepackage{hyperref}
\usepackage{enumitem}
\usepackage{fontspec}
\usepackage{amsmath}
\IfFontExistsTF{Microsoft YaHei}{\setCJKmainfont{Microsoft YaHei}}{\setCJKmainfont{SimSun}[BoldFont=SimHei]}
\IfFontExistsTF{Microsoft YaHei UI}{\setCJKsansfont{Microsoft YaHei UI}}{}
\IfFontExistsTF{Consolas}{\setmonofont{Consolas}}{}
\definecolor{linkblue}{RGB}{30,80,145}
\hypersetup{colorlinks=true,linkcolor=linkblue,urlcolor=linkblue,pdfborder={0 0 0}}
\setlength{\parindent}{2em}
\setlength{\parskip}{0.25em}
\linespread{1.16}
\setlist{itemsep=0.2em,topsep=0.35em,leftmargin=2.5em}
\pagestyle{plain}

\title{02\_渐进符号定义与证明}

\author{}

\date{}

\begin{document}

\maketitle

\section*{02 渐进符号定义与证明}

\subsection*{题目原文}

$O$、$\Omega$、$\Theta$ 的定义和相关函数证明

\subsection*{考查类型}

概念题和证明题。

\subsection*{三个定义}

设 $f(n)$ 和 $g(n)$ 是非负函数。

$f(n) = O(g(n))$ 表示 $f$ 被 $g$ 渐进上界控制：

\begin{quote}
\small\ttfamily\raggedright
\relax 存在正常数\ c\ 和\ n₀，使得对所有\ n\ ≥\ n₀，有\ 0\ ≤\ f(n)\ ≤\ c\ g(n)\\
\end{quote}

$f(n) = \Omega(g(n))$ 表示 $f$ 被 $g$ 渐进下界控制：

\begin{quote}
\small\ttfamily\raggedright
\relax 存在正常数\ c\ 和\ n₀，使得对所有\ n\ ≥\ n₀，有\ 0\ ≤\ c\ g(n)\ ≤\ f(n)\\
\end{quote}

$f(n) = \Theta(g(n))$ 表示二者同阶：

\begin{quote}
\small\ttfamily\raggedright
\relax 存在正常数\ c₁,\ c₂\ 和\ n₀，使得对所有\ n\ ≥\ n₀，有\ 0\ ≤\ c₁\ g(n)\ ≤\ f(n)\ ≤\ c₂\ g(n)\\
\end{quote}

等价地：

\begin{quote}
\small\ttfamily\raggedright
\relax f(n)\ =\ Θ(g(n))\ 当且仅当\ f(n)\ =\ O(g(n))\ 且\ f(n)\ =\ Ω(g(n))\\
\end{quote}

\subsection*{证明方法}

证明 $O$：找到一个常数 $c$，让 $f(n) \le c g(n)$ 在足够大的 $n$ 上成立。

证明 $\Omega$：找到一个常数 $c$，让 $f(n) \ge c g(n)$ 在足够大的 $n$ 上成立。

证明 $\Theta$：同时证明上界和下界。

\subsection*{示例}

证明 $3n^2 + 5n + 7 = O(n^2)$。

当 $n \ge 1$ 时：

\[3n^{2} + 5n + 7 \le  3n^{2} + 5n^{2} + 7n^{2} = 15n^{2}\]

取 $c = 15, n_0 = 1$，定义满足，因此结论成立。

证明 $3n^2 + 5n + 7 = \Omega(n^2)$。

当 $n \ge 1$ 时：

\[3n^{2} + 5n + 7 \ge  3n^{2}\]

取 $c = 3, n_0 = 1$，定义满足，因此结论成立。

由上界和下界同时成立，得到：

\[3n^{2} + 5n + 7 = \Theta(n^{2})\]

\subsection*{常用性质}

传递性：

\[f = O(g), g = O(h) \Rightarrow  f = O(h)\]

同阶对称性：

\[f = \Theta(g) \Rightarrow  g = \Theta(f)\]

多项式主项支配：

\[a_d n^{d} + ... + a_{1} n + a_{0} = \Theta(n^{d})\]

这里要求最高次项系数 $a_d > 0$。

\subsection*{例题}

\subsubsection*{例题 1}

题目：用定义证明：

\[5n^{3} + 2n^{2} + 7 = O(n^{3})\]

解答：

当 $n \ge 1$ 时：

\[5n^{3} + 2n^{2} + 7 \le  5n^{3} + 2n^{3} + 7n^{3} = 14n^{3}\]

取：

\[c = 14\]
\[n_{0} = 1\]

由 $O$ 的定义，结论成立。

\subsubsection*{例题 2}

题目：用定义证明：

\[5n^{3} + 2n^{2} + 7 = \Theta(n^{3})\]

解答：

上界由例题 1 得到：

\[5n^{3} + 2n^{2} + 7 \le  14n^{3}\]

下界为：

\[5n^{3} + 2n^{2} + 7 \ge  5n^{3}\]

取：

\[c_{1} = 5\]
\[c_{2} = 14\]
\[n_{0} = 1\]

满足：

\[c_{1} n^{3} \le  5n^{3} + 2n^{2} + 7 \le  c_{2} n^{3}\]

所以：

\[5n^{3} + 2n^{2} + 7 = \Theta(n^{3})\]

\subsection*{常见误区}

不要只代入几个数值说明渐进关系。渐进证明必须给出常数和阈值。

不要把 $O$ 当成等号意义上的相等。$O$ 表示上界关系，不表示精确增长。

\end{document}
